\subsection*{David POUJOL 1870-1939}

Sa famille « sortait » des Oubrets. Après ses études à Montauban, il fut pasteur à
Mazamet pendant de nombreuses années. Très libéral, il faisait preuve d’une
éloquence remarquable, attirant les foules. Cela expliquerait-il, en partie tout au
moins, la naissance d’un petit noyau évangélique rattaché aux Eglises Libres ?
Il perdit un de ses fils à la bataille de la Marne et, très patriote, comme on l’était à
cette époque, donna de nombreuses conférences dans toute la France dans lesquelles
il exaltait le courage de nos troupes face à la barbarie allemande. Ce qui ne laisse pas
de nous étonner aujourd’hui, car on peut penser que l’église dont il était le pasteur,
tout comme les responsables des Eglises Réformées lui laissaient toute liberté !
\emph{(voir son portrait dans « les souvenirs du pasteur Jean Gounelle »)}

\subsection*{Frank ROBERT \hspace{1em} 1909 - 1988}

Pasteur de Meyrueis de 1941 à 1946, Frank Robert fut un des responsables de la
résistance à Meyrueis ; il contribua au sauvetage de plusieurs familles juives , de
résistants réfugiés à Meyrueis et ses environs en particulier après l’anéantissement
du maquis Bir Hakeim à la Parade.
Membre du Comité de Libération, il fut adjoint au maire de 1944 et 1945.
Reconnu comme Juste des Nations ( Yad Vaschem à Jérusalem )

\subsection*{Rosine ROUEL 1881-1902}

Née à Gatuzieres. Après de très brillantes études elle fut nommée à l’âge de 20
ans, directrice du fameux cours complémentaire de filles de Vialas en Lozère.
Malheureusement elle était emportée un an plus tard, sans doute par la tuberculose.

\subsection*{Pierre-François SAMUEL, dit « François » \hspace{1em} 1745-1833}

Pasteur de Meyrueis, il est surtout connu par ses prises de position politiques au
moment de la Révolution. Il prend fait et cause pour les Girondins contre
Robespierre : il envoie même une adresse à la Convention contre la tyrannie de la
Montagne. Il essaya de faire passer la Lozère dans le camp des fédéralistes.
Emprisonné, il était en passe de finir sous la guillotine. Mais « ce vigoureux petit
pasteur n’attendit pas, il fit une corde avec ses draps et s’évada, 30 dragons à ses
trousses ».

\emph{(voir les pages : « C’était ça aussi la Révolution à Meyrueis ».}

Il trouva refuge dans le Tarn, avec sa femme et sa fille. Il s’efforça de se justifier
dans une longue lettre et ne fut plus jamais inquiété. Il a été pasteur à Millau et
Puylaurens, ce qui ne l’empêche pas de protester contre les règlements imposant
aux protestants de garnir leurs maisons de tentures lors de la Fête-Dieu, le 2\textsuperscript{o}
dimanche après Pentecôte.

\footnote{\emph{A la Fête-Dieu, l’hostie sainte, portée en grande pompe dans les rues, était proposée à l’adoration des fidèles catholiques. A Meyrueis, dans les années 1940, les maisons catholiques étaient ornées de draps tendus sur les façades, ainsi que de reposoirs avec des bouquets, des décorations, des images pieuses, pour le passage de la procession.}}
